\section{APÉNDICE B. LISTAS DE VERIFICACIÓN DE AUDITORÍA DE PROCESOS DE INGENIERÍA DE SOFTWARE GENERALES}

\begin{table}[H]
\centering
\caption{Lista de Verificación de Auditoría del Proceso de Planificación del Proyecto}
\label{fig:b1}
\begin{tabular}{|p{14cm}|}
\hline
\textbf{Proyecto:} \\
\textbf{Fecha:} \\
\textbf{Preparado por:} \\
\hline
\textbf{Procedimientos:} \\
\_\_\_\_ Los Planes y compromisos del Proyecto existen y están documentados en el SDP. \\
\_\_\_\_ Los estándares que rigen el proceso de desarrollo de software del proyecto están documentados y reflejados en el SDP. \\
\_\_\_\_ El contenido del SDP refleja la implementación consistente del proceso de software estándar de la organización. \\
\_\_\_\_ El SDP está bajo gestión de configuración. \\
\_\_\_\_ Las actividades de estimación de software se llevan a cabo de acuerdo con el Proceso de Estimación de Software y los resultados están documentados. \\
\_\_\_\_ La base de datos de la organización se utiliza para realizar estimaciones. \\
\_\_\_\_ Los requisitos de software son la base para los planes, productos de trabajo y actividades de software. \\
\_\_\_\_ Existen planes para llevar a cabo la gestión de configuración del software y están documentados en el SDP o en un Plan de Gestión de Configuración de Software (SCM Plan) separado. \\
\_\_\_\_ El Plan SCM está bajo gestión de configuración. \\
\_\_\_\_ Existen planes para llevar a cabo el aseguramiento de la calidad del software y están documentados en el SDP o en un Plan de Aseguramiento de la Calidad del Software (SQA Plan) separado. \\
\_\_\_\_ Existen planes para llevar a cabo las pruebas de integración del software y están documentados en un Plan de Pruebas de Software (STP). \\
\_\_\_\_ Existen planes para llevar a cabo las pruebas del sistema y están documentados en un STP. \\
\_\_\_\_ El STP está bajo gestión de configuración. \\
\_\_\_\_ Existen planes para llevar a cabo la transición del software y están documentados en un Plan de Transición de Software o en el Plan de Desarrollo de Software (SDP). \\
\_\_\_\_ Los cronogramas y planes del proyecto se someten a una revisión por pares antes de establecer sus líneas base. \\
\hline
\end{tabular}
\end{table}

\begin{table}[H]
\centering
\caption{Lista de Verificación de Auditoría del Proceso de Análisis de Requisitos del Sistema}
\label{fig:b3}
\begin{tabular}{|p{14cm}|}
\hline
\textbf{Proyecto:} \\
\textbf{Fecha:} \\
\textbf{Preparado por:} \\
\hline
\textbf{Procedimientos:} \\
\_\_\_\_ Los participantes correctos están involucrados en el proceso de análisis de requisitos del sistema para identificar todas las necesidades del usuario. \\
\_\_\_\_ Los requisitos se revisan para determinar si son factibles de implementar, están claramente establecidos y son consistentes. \\
\_\_\_\_ Los cambios en los requisitos asignados, los productos de trabajo y las actividades se identifican, revisan y rastrean hasta su cierre. \\
\_\_\_\_ El personal del proyecto involucrado en el proceso de análisis de requisitos del sistema está capacitado en los procedimientos y estándares necesarios aplicables a su área de responsabilidad para hacer el trabajo correctamente. \\
\_\_\_\_ Los compromisos resultantes de los requisitos asignados son negociados y acordados por los grupos afectados. \\
\_\_\_\_ Los compromisos están documentados, revisados, aceptados, aprobados y comunicados. \\
\_\_\_\_ Los requisitos asignados identificados como potencialmente problemáticos se revisan con el grupo responsable de analizar los requisitos y documentos del sistema, y se realizan los cambios necesarios. \\
\_\_\_\_ Se siguen y documentan los procesos prescritos para definir, documentar y asignar requisitos. \\
\_\_\_\_ Existe un proceso de CM para controlar y gestionar la línea base. \\
\_\_\_\_ Los requisitos están documentados, gestionados, controlados y rastreados (preferiblemente mediante una matriz). \\
\_\_\_\_ Los requisitos acordados se abordan en el SDP. \\
\hline
\end{tabular}
\end{table}

\begin{table}[H]
\centering
\caption{Lista de Verificación de Auditoría del Proceso de Análisis de Requisitos de Software}
\label{fig:b5}
\begin{tabular}{|p{14cm}|}
\hline
\textbf{Proyecto:} \\
\textbf{Fecha:} \\
\textbf{Preparado por:} \\
\hline
\textbf{Procedimientos:} \\
\textbf{Parte 1. Requisitos de Software} \\
\_\_\_\_ Los requisitos de software están documentados en una Especificación de Requisitos de Software (SRS) u otro formato aprobado y se incluyen en una matriz de trazabilidad. \\
\_\_\_\_ La SRS se mantiene bajo gestión de configuración. \\
\_\_\_\_ Los cambios en la SRS se someten a una revisión por pares antes de incorporarse a la línea base de requisitos. \\
\_\_\_\_ Asegurar que la revisión por pares valide la capacidad de prueba de los requisitos y que se establezcan métricas apropiadas para validar los requisitos de rendimiento medibles. \\
\_\_\_\_ Los planes, productos de trabajo y actividades de desarrollo de software se cambian para ser consistentes con los cambios en los requisitos de software. \\
\_\_\_\_ Las técnicas de análisis de requisitos de software son consistentes con el SDP. \\
\_\_\_\_ Las herramientas automatizadas adquiridas para gestionar los informes de problemas de requisitos de software y las solicitudes de cambio se utilizan correctamente. \\
\_\_\_\_ El grupo de ingeniería de software está capacitado para realizar actividades de gestión de requisitos. \\
\_\_\_\_ Se realizan mediciones y se utilizan para determinar el estado de la gestión de requisitos. \\
\textbf{Parte 2. Requisitos de Interfaz} \\
\_\_\_\_ Los requisitos de interfaz están documentados en una Especificación de Requisitos de Interfaz (IRS) u otro formato aprobado. \\
\_\_\_\_ La IRS se mantiene bajo gestión de configuración. \\
\_\_\_\_ Los cambios en la IRS se someten a una revisión por pares antes de incorporarse a la línea base de requisitos. \\
\_\_\_\_ Los planes, productos de trabajo y actividades de desarrollo de software se cambian para ser consistentes con los cambios en los requisitos de interfaz. \\
\_\_\_\_ Se utilizan herramientas automatizadas para gestionar los informes de problemas de requisitos de interfaz y las solicitudes de cambio. \\
\_\_\_\_ El grupo de ingeniería de software está capacitado para realizar actividades de gestión de requisitos. \\
\hline
\end{tabular}
\end{table}

\begin{table}[H]
\centering
\caption{Lista de Verificación de Auditoría del Proceso de Implementación y Pruebas Unitarias de Software}
\label{fig:b7}
\begin{tabular}{|p{14cm}|}
\hline
\textbf{Proyecto:} \\
\textbf{Fecha:} \\
\textbf{Preparado por:} \\
\hline
\textbf{Procedimientos:} \\
\textbf{Parte 1. Implementación de Software} \\
\_\_\_\_ El código y la matriz de trazabilidad se preparan y se mantienen actualizados y consistentes en función de los cambios aprobados en los requisitos de software. \\
\_\_\_\_ Las revisiones del código (revisión por pares) evalúan el cumplimiento del código con el diseño aprobado, identifican defectos en el código y se evalúan y reportan alternativas. \\
\_\_\_\_ Las revisiones del código se llevan a cabo de acuerdo con el Proceso de Revisión por Pares. \\
\_\_\_\_ Los cambios en el código se identifican, revisan y rastrean hasta su cierre. \\
\_\_\_\_ El código se mantiene bajo gestión de configuración. \\
\_\_\_\_ Los cambios en el código se someten a una revisión por pares antes de incorporarse a la línea base de software. \\
\_\_\_\_ La codificación del software es consistente con la metodología de codificación aprobada en el Plan de Desarrollo de Software (SDP). \\
\_\_\_\_ El método, como la Carpeta de Desarrollo de Software o la Carpeta de Desarrollo de Unidades, utilizado para rastrear y documentar el desarrollo/mantenimiento de una unidad de software está implementado y se mantiene actualizado. \\
\textbf{Parte 2. Pruebas Unitarias} \\
\_\_\_\_ Las pruebas unitarias de software se llevan a cabo de acuerdo con los estándares y procedimientos aprobados descritos en el SDP. \\
\_\_\_\_ Asegurar que los criterios de aprobación para las pruebas unitarias estén documentados y que se haya registrado el cumplimiento. \\
\_\_\_\_ Los resultados de las pruebas unitarias se documentan en la Carpeta de Desarrollo de Software o en la Carpeta de Desarrollo de Unidades. \\
\hline
\end{tabular}
\end{table}

\begin{table}[H]
\centering
\caption{Lista de Verificación de Auditoría del Proceso de Integración y Pruebas Unitarias}
\label{fig:b8}
\begin{tabular}{|p{14cm}|}
\hline
\textbf{Proyecto:} \\
\textbf{Fecha:} \\
\textbf{Preparado por:} \\
\hline
\textbf{Procedimientos:} \\
\textbf{Parte 1. Identificación de Configuración} \\
\_\_\_\_ Los CI se integran bajo control de cambios. Si no, vaya a la Parte 2; de lo contrario, continúe. \\
\_\_\_\_ Los CI integrados en el sistema se obtuvieron de un representante autorizado de Gestión de Configuración de acuerdo con el SDP. \\
\_\_\_\_ Las versiones base de cada CI se integran en el sistema. \\
\_\_\_\_ Todos los componentes de CI de la integración de software están bajo control de cambios de acuerdo con el Plan SCM. \\
\_\_\_\_ Si es así, describa cómo se identifican. \\
\textbf{Parte 2. Proceso de Integración} \\
\_\_\_\_ Existe un plan para la integración de los CI. \\
\_\_\_\_ El plan especifica el orden y el cronograma en el que se integran los CI. \\
\_\_\_\_ Los CI se integran de acuerdo con el cronograma y en el orden especificado. \\
\_\_\_\_ El plan de integración especifica qué versión de cada CSC y CSU se va a integrar. \\
\_\_\_\_ Se integra la versión correcta. \\
\_\_\_\_ Los CI integrados han completado las pruebas unitarias. \\
\_\_\_\_ Se han completado las correcciones requeridas. \\
\_\_\_\_ Los CI han sido reprobados. \\
\_\_\_\_ Los procedimientos de prueba están definidos para la integración de CI. \\
\_\_\_\_ Se siguen los procedimientos definidos. \\
\_\_\_\_ Los casos de prueba están definidos. \\
\_\_\_\_ Se siguen los casos de prueba definidos. \\
\_\_\_\_ Los criterios de aprobación/reprobación de las pruebas están definidos. \\
\_\_\_\_ Se siguen los criterios de aprobación/reprobación definidos. \\
\_\_\_\_ Los resultados de las pruebas se documentan en las Carpetas de Desarrollo de Unidades. \\
\hline
\end{tabular}
\end{table}

\begin{table}[H]
\centering
\caption{Lista de Verificación de Auditoría del Proceso de Pruebas de Calificación de CI}
\label{fig:b9}
\begin{tabular}{|p{14cm}|}
\hline
\textbf{Proyecto:} \\
\textbf{Fecha:} \\
\textbf{Preparado por:} \\
\hline
\textbf{Procedimientos:} \\
\textbf{Parte 1. Plan de Pruebas} \\
\_\_\_\_ Existe un plan de pruebas y descripciones de pruebas aprobados. \\
\_\_\_\_ El plan y las descripciones utilizados están bajo control de gestión de configuración. \\
\_\_\_\_ Se utiliza la última versión del plan y la descripción. \\
\textbf{Parte 2. Proceso de Pruebas} \\
\_\_\_\_ El software del sistema se recibió de una fuente autorizada de gestión de configuración. \\
\_\_\_\_ El entorno de prueba, incluidos los requisitos de hardware y software, se configura según lo requiere el plan de pruebas. \\
\_\_\_\_ Las pruebas se realizaron en el orden correcto. \\
\_\_\_\_ Se ejecutó cada caso de prueba en la descripción de la prueba. \\
\_\_\_\_ El sistema se prueba después de cada integración de CI. \\
\_\_\_\_ Los criterios de aprobación de las pruebas están documentados y se registra el cumplimiento. \\
\_\_\_\_ Los resultados de las pruebas se registran en un informe de pruebas. \\
\_\_\_\_ Si es así, describa la información que se registra y dónde se registra. \\
\_\_\_\_ Los CI se vuelven a probar después de la integración para asegurar que aún cumplen con sus requisitos sin interferencia del resto del sistema. \\
\_\_\_\_ El sistema que resulta de la integración de cada CI se coloca bajo control de gestión de configuración. \\
\_\_\_\_ Si es así, describa cómo se identifica. \\
\textbf{Parte 3. Reporte de Problemas} \\
\_\_\_\_ Las discrepancias encontradas se ingresan en un sistema de gestión de configuración P/CR o STR para el control de cambios. Si no, describa dónde se registran. \\
\_\_\_\_ Las entradas se completan en el momento en que se encuentran las discrepancias. \\
\_\_\_\_ El número de referencia del P/CR o STR se mantiene en el archivo de pruebas. Si no, describa dónde se guarda. \\
\_\_\_\_ Los P/CR o STR se escriben cuando se encuentran problemas en el entorno de prueba, el plan de pruebas, las descripciones de las pruebas o los casos de prueba. \\
\_\_\_\_ Estos P/CR o STR se envían a través del mismo proceso de control de cambios que las discrepancias de software. \\
\textbf{Parte 4. Modificaciones} \\
\_\_\_\_ ¿Se realizan modificaciones o correcciones en el entorno de prueba durante las pruebas? \\
\_\_\_\_ Si es así, ¿se aprobaron las modificaciones a través del proceso de control de cambios antes de la implementación? \\
\_\_\_\_ ¿Qué documentación de las modificaciones existe? \\
\_\_\_\_ ¿Cuáles son los cambios? \\
\_\_\_\_ ¿Se realizan modificaciones o correcciones en las descripciones de las pruebas durante las pruebas? \\
\_\_\_\_ Si es así, ¿se aprobaron las modificaciones por el proceso de control de cambios antes de la implementación? \\
\_\_\_\_ ¿Qué documentación de las modificaciones existe? \\
\_\_\_\_ ¿Cuál es la fecha de aprobación del LCCB? \\
\_\_\_\_ ¿Cuál es la fecha de lanzamiento del cambio STD? \\
\_\_\_\_ ¿Se siguieron los procedimientos de control de cambios en el Plan SCM? \\
\_\_\_\_ Si no, ¿cuáles son los cambios? \\
\textbf{Parte 5. Criterios de Finalización} \\
\_\_\_\_ Todos los CI especificados se han integrado en el sistema. \\
\_\_\_\_ Todos los casos de prueba se han ejecutado en el sistema. \\
\_\_\_\_ Todos los P/CR o STR están cerrados. \\
\_\_\_\_ Si no, todos los P/CR o STR pendientes están documentados adecuadamente en el VDD. \\
\_\_\_\_ El Informe de Pruebas de Software se ha completado y aprobado. \\
\_\_\_\_ El Informe de Pruebas de Software se ha colocado bajo control de cambios. \\
\_\_\_\_ La autoridad competente ha determinado si el sistema pasó o falló las pruebas de integración. \\
\_\_\_\_ Indique el individuo o grupo que determinó si el sistema pasó o falló. \\
\_\_\_\_ Describa cómo se realizó la determinación de aprobación o fallo. \\
\_\_\_\_ El sistema de software está listo para ser integrado con el sistema operativo. \\
\textbf{Parte 6. Archivos de Desarrollo de Software} \\
\_\_\_\_ El Informe de Pruebas de Software incluye cualquier repetición de pruebas debido a fallos de software. \\
\_\_\_\_ Si es así, enumere los fallos con sus correspondientes números de referencia P/CR o STR. \\
\_\_\_\_ Utilizando el sistema de CM P/CR o STR, enumere todos los CI cambiados debido a estos fallos. \\
\_\_\_\_ Todos los archivos de desarrollo de software de los CI enumerados se actualizaron de acuerdo con el SDP. \\
\_\_\_\_ Si no, enumere todos los archivos de desarrollo de software que no se actualizaron. \\
\textbf{Parte 7. Precisión del Informe de Pruebas de Software} \\
\_\_\_\_ El Informe de Pruebas de Software proporciona el Número de Identificación de Configuración (CIN) para todos los documentos de prueba (STP, STD) y software. Si no, el Informe de Pruebas de Software está incompleto. \\
\_\_\_\_ ¿El probador puede ejecutar las pruebas de evaluación con los CIN especificados? Si no, el Informe de Pruebas de Software es inexacto. \\
\_\_\_\_ ¿Los resultados de estas pruebas coinciden con el Informe de Pruebas de Software? Si no, el Informe de Pruebas de Software es inexacto. \\
\hline
\end{tabular}
\end{table}

\begin{table}[H]
\centering
\caption{Lista de Verificación de Auditoría del Proceso de Entrega de Productos Finales}
\label{fig:b10}
\begin{tabular}{|p{14cm}|}
\hline
\textbf{Proyecto:} \\
\textbf{Fecha:} \\
\textbf{Preparado por:} \\
\hline
\textbf{Procedimientos:} \\
\textbf{Parte 1. Auditorías} \\
\_\_\_\_ Se realizó una auditoría funcional, si era requerida. \\
\_\_\_\_ Se realizó una auditoría física, si era requerida. \\
\textbf{Parte 2. Generación del Paquete} \\
\_\_\_\_ El software se genera a partir de la biblioteca de software de acuerdo con el SDP. \\
\_\_\_\_ Si es así, el software es la última versión del software en la biblioteca de software. \\
\_\_\_\_ Si no, ¿por qué no? \\
\_\_\_\_ La documentación se genera a partir de los originales controlados por el personal de Gestión de Configuración según lo requiere el Plan SCM. \\
\textbf{Parte 3. Paquete de Entrega} \\
\_\_\_\_ El soporte del software está etiquetado correctamente, mostrando como mínimo el nombre del software, la fecha de lanzamiento y el número de versión correcto. La lista de verificación en la Figura B-12 contiene más detalles sobre el etiquetado de soportes. \\
\_\_\_\_ El Documento de Descripción de la Versión está con el soporte del software. \\
\_\_\_\_ Si es así, el Documento de Descripción de la Versión es la versión correcta para el software. \\
\_\_\_\_ El Documento de Descripción de la Versión ha sido inspeccionado formalmente. \\
\_\_\_\_ El Manual del Usuario está con el soporte del software. \\
\_\_\_\_ Si es así, el Manual del Usuario es la versión correcta para el software. \\
\_\_\_\_ El Manual del Usuario ha sido inspeccionado formalmente. \\
\textbf{Parte 4. Lista de Distribución de Soportes} \\
\_\_\_\_ Existe una lista de distribución para el entregable. \\
\_\_\_\_ Si es así, está completa, todas las organizaciones enumeradas, todas las direcciones correctas y actuales. \\
\_\_\_\_ Si es así, ¿hay alguna organización enumerada que no necesite recibir el entregable? \\
\_\_\_\_ ¿El entregable está clasificado? \\
\_\_\_\_ Si es así, el personal en la lista de distribución tiene la autorización y necesidad de conocer requeridas. \\
\textbf{Parte 5. Empaquetado} \\
\_\_\_\_ El material de empaquetado es adecuado para el contenido y el método de transmisión utilizado. \\
\_\_\_\_ ¿El paquete contiene una carta de transmisión firmada? \\
\_\_\_\_ Si es así, la información de la transmisión es correcta. \\
\_\_\_\_ Todo el contenido está listado en la transmisión contenida en el paquete. \\
\_\_\_\_ ¿El paquete incluye un formulario de acuse de recibo? \\
\textbf{Parte 6. Notificación de Problemas} \\
\_\_\_\_ Existe un método especificado para que la organización receptora notifique al remitente sobre problemas y deficiencias en el paquete. \\
\_\_\_\_ Si es así, describa el método. \\
\_\_\_\_ Existe un método especificado para registrar y manejar los problemas de distribución. Descríbalo. \\
\_\_\_\_ Los problemas de distribución son manejados por una persona específica. Identifique a esa persona. \\
\hline
\end{tabular}
\end{table}

\begin{table}[H]
\centering
\caption{Lista de Verificación de Auditoría del Proceso de Certificación de Soportes}
\label{fig:b12}
\begin{tabular}{|p{14cm}|}
\hline
\textbf{Proyecto:} \\
\textbf{Fecha:} \\
\textbf{Preparado por:} \\
\hline
\textbf{Procedimientos:} \\
\textbf{Parte 1. Producción de Soportes} \\
\_\_\_\_ El soporte que contiene el código fuente y el soporte que contiene el código objeto que se entregan se corresponden entre sí. \\
\_\_\_\_ Existe un proceso documentado que se utiliza para implementar la producción de soportes a partir de la biblioteca de software. \\
\_\_\_\_ Si no, se está creando un plan o se está documentando una razón por la que no se necesita uno. Vaya a la Parte 2. \\
\_\_\_\_ El plan se siguió en la producción de soportes. \\
\_\_\_\_ El software se creó a partir de los archivos correctos en la biblioteca de software por personal de CM. \\
\_\_\_\_ Los documentos se crearon a partir de copias maestras aprobadas por personal de CM. \\
\textbf{Parte 2. Etiquetado de Soportes} \\
\_\_\_\_ Existe un estándar documentado que se sigue en el etiquetado de los soportes. \\
\_\_\_\_ Si es así, describa el método estándar utilizado para identificar el producto, la versión y el Número de Identificación de Configuración. \\
\_\_\_\_ El soporte está claramente etiquetado. \\
\_\_\_\_ La etiqueta contiene toda la información requerida (producto, versión y Número de Identificación de Configuración). \\
\_\_\_\_ Si el software está clasificado, el soporte refleja claramente la clasificación correcta. \\
\_\_\_\_ El documento de software está claramente etiquetado con el producto, CIN y número de versión, si corresponde. \\
\textbf{Parte 3. Contenidos del Soporte} \\
\_\_\_\_ Existe un listado del contenido en el soporte. \\
\_\_\_\_ Si es así, describa dónde se encuentra el listado. \\
\_\_\_\_ El soporte contiene el contenido especificado en el listado. \\
\_\_\_\_ El contenido del soporte coincide con la información de la etiqueta, es decir, ¿es la versión correcta para la plataforma de hardware correcta? \\
\_\_\_\_ Los documentos contienen todas las páginas de cambio requeridas para esta versión de los documentos. \\
\hline
\end{tabular}
\end{table}

\begin{table}[H]
\centering
\caption{Lista de Verificación de Auditoría del Proceso de Certificación de Software No Entregable}
\label{fig:b13}
\begin{tabular}{|p{14cm}|}
\hline
\textbf{Proyecto:} \\
\textbf{Fecha:} \\
\textbf{Preparado por:} \\
\hline
\textbf{Procedimientos:} \\
\_\_\_\_ El software entregable depende del software no entregable. \\
\_\_\_\_ Si es así, se ha dispuesto que el adquiriente tenga o pueda obtener el software no entregable. \\
\_\_\_\_ El software no entregado cumple con el uso previsto. \\
\_\_\_\_ El software no entregado se coloca bajo gestión de configuración. \\
\hline
\end{tabular}
\end{table}

\begin{table}[H]
\centering
\caption{Lista de Verificación de Auditoría del Proceso de Almacenamiento y Manipulación}
\label{fig:b14}
\begin{tabular}{|p{14cm}|}
\hline
\textbf{Proyecto:} \\
\textbf{Fecha:} \\
\textbf{Preparado por:} \\
\hline
\textbf{Procedimientos:} \\
\_\_\_\_ Los documentos y soportes se almacenan de acuerdo con el procedimiento de la Biblioteca de Desarrollo de Software. \\
\_\_\_\_ Las áreas de almacenamiento de productos en papel están libres de efectos ambientales adversos (alta humedad, fuerzas magnéticas, calor y polvo). \\
\_\_\_\_ Las áreas de almacenamiento de productos en soportes están libres de efectos ambientales adversos (alta humedad, fuerzas magnéticas, calor y polvo). \\
\_\_\_\_ Los contenedores de almacenamiento para material clasificado son apropiados para el nivel de material clasificado. \\
\hline
\end{tabular}
\end{table}

\begin{table}[H]
\centering
\caption{Lista de Verificación de Auditoría del Proceso de Desviaciones y Exenciones}
\label{fig:b16}
\begin{tabular}{|p{14cm}|}
\hline
\textbf{Proyecto:} \\
\textbf{Fecha:} \\
\textbf{Preparado por:} \\
\hline
\textbf{Procedimientos:} \\
\_\_\_\_ Las desviaciones y exenciones se preparan de acuerdo con los procedimientos documentados en el Plan SCM del proyecto. \\
\_\_\_\_ Las desviaciones y exenciones son revisadas y aprobadas por el personal apropiado de acuerdo con el Plan SCM del proyecto. \\
\_\_\_\_ Los registros de desviaciones y exenciones se mantienen y reflejan la configuración de desarrollo actual. \\
\hline
\end{tabular}
\end{table}

\begin{table}[H]
\centering
\caption{Lista de Verificación de Auditoría del Proceso de Gestión de Configuración de Software}
\label{fig:b17}
\begin{tabular}{|p{14cm}|}
\hline
\textbf{Proyecto:} \\
\textbf{Fecha:} \\
\textbf{Preparado por:} \\
\hline
\textbf{Procedimientos:} \\
\textbf{Parte 1. Plan SCM} \\
\_\_\_\_ El proyecto sigue la política organizacional para implementar SCM. \\
\_\_\_\_ Existe el grupo responsable de coordinar e implementar SCM para el proyecto. \\
\_\_\_\_ Se utiliza un Plan SCM documentado y aprobado como base para realizar las actividades de SCM. \\
\_\_\_\_ El control de configuración de los cambios en los documentos y software de línea base se gestiona de acuerdo con el Plan SCM. \\
\_\_\_\_ Se establece un sistema de biblioteca de gestión de configuración como repositorio de la línea base de software. \\
\_\_\_\_ La biblioteca de CM es el único lugar de almacenamiento de la versión base de todo el software. \\
\_\_\_\_ El acceso a los productos de software en la biblioteca de CM se realiza de acuerdo con los procedimientos de Control de la Biblioteca. \\
\_\_\_\_ Los productos de trabajo de software que se colocarán bajo SCM se identifican de acuerdo con el Plan SCM. \\
\_\_\_\_ Existe la Junta de Control de Cambios Local (LCCB) e implementa los procedimientos de la LCCB. \\
\_\_\_\_ Las solicitudes de cambio y los informes de problemas para todos los elementos de configuración se manejan de acuerdo con el procedimiento de PCR. \\
\_\_\_\_ Los cambios en las líneas base se controlan de acuerdo con el Plan SCM, el procedimiento de la LCCB y el procedimiento de PCR. \\
\_\_\_\_ Los productos de la biblioteca de línea base de software se crean y su liberación se controla de acuerdo con los procedimientos de Control de la Biblioteca. \\
\_\_\_\_ Los informes de contabilidad del estado de configuración se preparan de acuerdo con el Plan SCM. \\
\textbf{Parte 2. Identificación de Configuración} \\
\_\_\_\_ Se pueden identificar las Líneas Base de Productos y la Biblioteca de Desarrollo. \\
\_\_\_\_ Si es así, describa el método utilizado para identificar las Líneas Base y la Biblioteca de Desarrollo. \\
\_\_\_\_ Enumere los documentos que componen estas Líneas Base y la Biblioteca de Desarrollo. \\
\_\_\_\_ Se puede identificar la documentación y los soportes de almacenamiento informático que contienen código, documentación o ambos. \\
\_\_\_\_ Si es así, describa el método utilizado para identificar la documentación y los soportes de almacenamiento informático y enumere los documentos que se colocan bajo control de configuración. \\
\_\_\_\_ Se puede identificar cada CI y sus componentes correspondientes. \\
\_\_\_\_ Si es así, describa el método utilizado para identificarlos. \\
\_\_\_\_ Se utiliza un método para identificar el nombre, la versión, el lanzamiento, el estado del cambio y cualquier otro detalle de identificación de cada elemento entregable. \\
\_\_\_\_ Si es así, describa el método utilizado. \\
\_\_\_\_ Para cada cliente, identifique el elemento entregable, la versión, el lanzamiento y el estado del cambio que se está utilizando. \\
\_\_\_\_ Se utiliza un método para identificar la versión de cada CI a la que se aplica la documentación de software correspondiente. \\
\_\_\_\_ Si es así, describa el método utilizado. \\
\_\_\_\_ Enumere la SRS y el SDD para cada CI. \\
\_\_\_\_ Se utiliza un método para identificar la versión específica del software contenido en un soporte entregable, incluidos todos los cambios incorporados desde su lanzamiento anterior. \\
\_\_\_\_ Si es así, describa el método utilizado. \\
\_\_\_\_ ¿El soporte entregable coincide con los originales de CM? Enumere cualquier discrepancia. \\
\textbf{Parte 3. Control de Configuración} \\
\_\_\_\_ Existe un plan establecido para realizar el control de configuración. \\
\_\_\_\_ Si es así, existe un método para establecer una Biblioteca de Desarrollo para cada CI. \\
\_\_\_\_ Si es así, describa el método utilizado. \\
\_\_\_\_ Enumere los CI en la Biblioteca de Desarrollo. \\
\_\_\_\_ Existe un método para mantener copias actualizadas de la documentación y el código entregable. \\
\_\_\_\_ Si es así, describa el método utilizado. \\
\_\_\_\_ Enumere las copias actuales. Enumere todas las discrepancias. \\
\_\_\_\_ Existe un método para controlar la preparación y difusión de cambios en las copias maestras del software y la documentación entregable. \\
\_\_\_\_ Si es así, describa el método utilizado. \\
\_\_\_\_ Las copias maestras de los entregables reflejan solo los cambios aprobados. Enumere cualquier discrepancia. \\
\_\_\_\_ Enumere los cambios en el software/documentos entregables actuales. \\
\textbf{Parte 4. Contabilidad del Estado de Configuración} \\
\_\_\_\_ Existe un plan documentado para implementar y realizar la contabilidad del estado de configuración. \\
\_\_\_\_ Existen informes de estado sobre todos los productos que componen las Bibliotecas de Desarrollo y las Líneas Base Funcionales Asignadas y de Producto. \\
\_\_\_\_ Los cambios propuestos e implementados en un CI y sus documentos de identificación de configuración asociados se registran y reportan. \\
\_\_\_\_ Si es afirmativo en dos de cada tres, responda lo siguiente. Si no, vaya a la Parte 5. \\
\_\_\_\_ Describa el método utilizado para proporcionar trazabilidad de los cambios a los productos controlados. \\
\_\_\_\_ Describa el método utilizado para comunicar el estado de la identificación de configuración y el software asociado. \\
\_\_\_\_ Describa el método para asegurar que los documentos entregados describan y representen el software asociado. \\
\textbf{Parte 5. Propuestas de Cambio de Ingeniería} \\
\_\_\_\_ Las Propuestas de Cambio de Ingeniería (ECP) se preparan de acuerdo con el estándar apropiado. \\
\_\_\_\_ Los Avisos de Cambio de Software (SCN) se preparan de acuerdo con el estándar apropiado. \\
\_\_\_\_ Describa el método utilizado para manejar los cambios solicitados al CI. \\
\_\_\_\_ Describa el método utilizado para autorizar SCN y ECP. \\
\hline
\end{tabular}
\end{table}

\begin{table}[H]
\centering
\caption{Lista de Verificación de Auditoría del Proceso de Control de la Biblioteca de Desarrollo de Software}
\label{fig:b18}
\begin{tabular}{|p{14cm}|}
\hline
\textbf{Proyecto:} \\
\textbf{Fecha:} \\
\textbf{Preparado por:} \\
\hline
\textbf{Procedimientos:} \\
\textbf{Parte 1. Establecimiento del Entorno y Control de la SDL} \\
\_\_\_\_ La SDL está establecida y existen procedimientos para gobernar su operación. \\
\_\_\_\_ La SDL proporciona un medio positivo para reconocer los elementos relacionados (es decir, aquellas versiones que constituyen una línea base particular y proteger el software contra la destrucción o modificación no autorizada). \\
\_\_\_\_ La documentación y los materiales del programa informático aprobados por la LCCB se colocan bajo control de la biblioteca. \\
\_\_\_\_ Todo el software, herramientas y documentación relevantes para el desarrollo de software se colocan bajo control de la biblioteca. \\
\_\_\_\_ Existen procedimientos/estándares publicados para la SDL. \\
\_\_\_\_ Los procedimientos de la SDL incluyen la identificación de las personas/organización responsables de recibir, almacenar, controlar y difundir los materiales de la biblioteca. \\
\_\_\_\_ El acceso a la SDL está limitado al personal autorizado. \\
\_\_\_\_ Si es así, describa los procedimientos utilizados para limitar el acceso. \\
\_\_\_\_ Existen salvaguardas para asegurar que no se realicen alteraciones no autorizadas al material controlado. \\
\_\_\_\_ Si es así, describa esas salvaguardas. \\
\_\_\_\_ Describa o enumere los materiales a controlar. \\
\_\_\_\_ Describa cómo se aprueban y colocan bajo control los materiales. \\
\_\_\_\_ Describa cómo se manejan los cambios en la parte del software y cómo se identifican los cambios en las líneas de código. \\
\_\_\_\_ La SDL contiene copias maestras de cada CI bajo control de la Biblioteca de Programas Informáticos. \\
\_\_\_\_ Se realiza una copia de seguridad periódica del software para evitar la pérdida de información debido a cualquier falla del Sistema de Biblioteca de Desarrollo. \\
\_\_\_\_ Si es así, describa el procedimiento de copia de seguridad y la frecuencia de las copias de seguridad. \\
\textbf{Parte 2. Garantía de Validez del Material Controlado} \\
\_\_\_\_ Las duplicaciones de copias maestras controladas y probadas se verifican antes de la entrega como copias exactas. \\
\_\_\_\_ Todos los productos de software entregables que se duplican de copias maestras controladas y probadas se comparan con esa copia maestra para asegurar la duplicación exacta. \\
\_\_\_\_ Describa cómo se asignan los números de identificación y los códigos de revisión a los documentos y software controlados. \\
\_\_\_\_ Describa cómo se registran las liberaciones de materiales controlados. \\
\_\_\_\_ Enumere las personas u organización responsables de la garantía de la validez del soporte de software. \\
\_\_\_\_ Existe un procedimiento de liberación formal. Si es así, ¿descríbalo? \\
\_\_\_\_ El material contenido en la biblioteca se actualiza de manera oportuna y correcta cuando se autoriza un cambio en cualquiera de estos materiales. \\
\hline
\end{tabular}
\end{table}

\begin{table}[H]
\centering
\caption{Lista de Verificación de Auditoría del Proceso de Software No Desarrollado}
\label{fig:b19}
\begin{tabular}{|p{14cm}|}
\hline
\textbf{Proyecto:} \\
\textbf{Fecha:} \\
\textbf{Preparado por:} \\
\hline
\textbf{Procedimientos:} \\
\_\_\_\_ El software no desarrollado cumple con las funciones previstas. \\
\_\_\_\_ El software no desarrollado se coloca bajo CM interno. \\
\_\_\_\_ Las disposiciones de derechos de datos y licencias son consistentes con el SDP. \\
\hline
\end{tabular}
\end{table}
